\documentclass[lettersize,journal]{IEEEtran}

\usepackage{amsmath,amsfonts,amssymb}
\usepackage{algorithmic}
\usepackage{algorithm}
\usepackage{array}
\usepackage{textcomp}
\usepackage{stfloats}
\usepackage{url}
\usepackage{cite}
\usepackage{graphicx}

\renewcommand{\algorithmicrequire}{\textbf{Input:}}
\renewcommand{\algorithmicensure}{\textbf{Output:}}
% Chinese text: use TeX Live bundled Fandol fonts for maximum portability.
% Compile with XeLaTeX. This avoids depending on Windows font names such as SimSun.
\usepackage[UTF8,fontset=fandol]{ctex}

\hyphenation{op-tical net-works semi-conduc-tor IEEE-Xplore}

\begin{document}
\title{Communication-Localization Fingerprint Aided Multi-UAV Trajectory Planning in Urban Environments}

\author{}

\maketitle

\begin{abstract}
城市低空环境的建筑遮挡会使得无人机与地面基站之间的通信链路质量和定位可观测性在空间上呈现显著的非连续性，即使采用通定一体化等新型技术手段，也难以完全避免通信中断或定位失锁等性能凹点。在此背景下，本文提出一种面向城市环境通信与定位服务的通信-定位指纹（Communication-Localization Fingerprint, CLF）辅助多无人机协同路径规划方法。首先，通过离线构建能够表征空间服务质量的通信-定位指纹地图，实现通信与定位性能的统一表示。随后，将多无人机协同路径规划问题建模为部分可观测马尔可夫决策过程，并采用集中训练、分布执行的多智能体强化学习框架进行求解。仿真结果表明，所提方法能够有效规避通信与定位性能退化区域，提高多无人机任务执行过程中的服务连续性和路径规划可靠性。
\end{abstract}

\begin{IEEEkeywords}
多无人机；通信—定位指纹；城市低空环境；路径规划；通信与定位一体化；多智能体强化学习
\end{IEEEkeywords}

\section{Introduction}

随着第六代 (6G) 无线通信系统的发展，通定一体化技术（ILAC）因其降低成本\cite{ref1}以及提高频谱效率等特点而广受关注。凭借视距链路、高机动性和快速部署等优势\cite{ref2,ref3,ref4,ref5,ref6}，无人机已成为城市低空经济中重要移动终端。近年来，大量研究聚焦于无人机辅助一体化系统的资源分配、路径规划、能源效率等方面。例如文献\cite{ref7}针对\textbf{城市环境中 UAV中继部署与三维位置优化问题}，提出一种\textbf{低复杂度空间搜索算法}以寻找最优UAV部署位置；针对\textbf{空地协同ILAC网络中UAV部署与接入控制问题}，研究\cite{ref8}提出基于\textbf{CRITIC指标评估的迭代贪婪搜索算法}以及\textbf{概率搜索分配方法}，用于优化UAV部署及通信与定位任务分配。

上述研究涉及的凸优化或启发式算法在静态环境中表现良好，然而在动态环境中，传统算法的计算复杂度较大、可行性较差\cite{ref9}。相比之下，强化学习可通过环境交互来改进策略，以减少决策过程中的复杂性\cite{ref10}。已有很多研究借助强化学习对无人机进行了资源分配与路径规划：例如，Wang等人在文献\cite{ref9}中构建了一种\textbf{分层多智能体深度强化学习（MADRL）框架}，联合优化 UAV-ISAC 网络中的用户关联、轨迹规划与功率分配问题；针对\textbf{山区环境中地形遮挡导致通信定位性能下降的问题}，文献\cite{ref11}构建了一个结合地理信息的强化学习决策框架 GeoAgg-HSAC以优化无人机轨迹与资源分配。该研究利用\textbf{图神经网络对受地形遮挡影响的网络状态进行聚合表示}，并通过 \textbf{HSAC 算法}进行混合动作空间决策；此外，Gao等人在文献\cite{ref12}中研究了\textbf{多无人机服务移动用户的UAV-ISAC场景}，提出一种\textbf{HR-MAPPO多智能体强化学习算法}，用于联合优化无人机轨迹与波束赋形。

尽管现有的研究方法能很好地完成无人机一体化系统的资源分配和路径规划，然而，很多研究通常需要在轨迹规划或资源分配过程中反复评估通信链路状态、定位几何结构以及误差传播对系统性能的影响。当城市环境规模增大、多无人机数量增加或约束条件复杂化时，这种在线评估机制将显著增加决策复杂度，不利于实时规划与协同决策。虽然已有某些基于三维地图无人机城市环境轨迹规划达到了很好的要求（\cite{ref13,ref14,ref15}），但是在考虑多无人机航路挤占、通信定位一体化、能量等多重约束下，原有的地图构建算法很难满足上述复杂问题的需求，整个系统在实际运转时可能面临定位失锁、通信断链和无人机碰撞的风险。

针对上述问题，本文提出一种基于通信-定位指纹（Communication-Localization Fingerprint, CLF）的多无人机协同路径规划框架。通过离线融合通信和定位性能，构建能够表征空间服务质量的通信-定位指纹地图，进而使得无人机能够在在线规划过程中快速识别通信与定位性能退化区域，避免重复进行复杂的通信与定位性能评估。在此基础上，构建多约束的多无人机协同路径规划问题，并将其建模为去中心化部分可观测马尔可夫决策过程（Dec-POMDP）。随后，在集中训练、分布执行（Centralized Training and Decentralized Execution, CTDE）的多智能体强化学习框架进行求解，实现多重约束下的长期稳定协同优化。

\section{System Model and Problem Formulation}

\subsection{Urban Environment Model}

考虑一个典型城市低空环境，区域内包含建筑物和地面基站。建筑物集合记为
$\mathbf{B}=\{b_1,b_2,\ldots,b_{|\mathbf{B}|}\}$，其中第 $j$ 个建筑物 $b_j$ 被建模为与坐标轴平行的长方体，其中心坐标和尺寸分别为
$\mathbf{c}_j^b=(x_j^b,y_j^b,z_j^b)$ 和 $(l_j,w_j,h_j)$。地面基站集合记为
$\mathbf{M}=\{m_1,m_2,\ldots,m_{|\mathbf{M}|}\}$，其中第 $k$ 个基站 $m_k$ 的位置、发射功率和载波频率分别为
$\mathbf{p}_k^m=(x_k^m,y_k^m,z_k^m)$、$P_k$ 和 $f_c$。

为便于路径规划与环境建模，本文将三维空域离散为体素集合
$\mathbf{V}=\{v_1,v_2,\ldots,v_{|\mathbf{V}|}\}$。第 $i$ 个体素 $v_i$ 的中心坐标为
$\mathbf{c}_i^v=(x_i^v,y_i^v,z_i^v)$，边长为 $\Delta r$，其占用区域记为 $\Omega(v_i)$。建筑物 $b_j$ 的占用区域记为 $\Omega(b_j)$。体素是否被建筑物占用由指示函数刻画：
\begin{equation}
I_{\mathrm{occ}}(v_i)=
\begin{cases}
1, & \exists b_j\in\mathbf{B},\ \Omega(v_i)\cap\Omega(b_j)\neq\varnothing,\\
0, & \mathrm{otherwise}
\end{cases}
\end{equation}

考虑建筑物避障、安全距离和低空飞行高度约束，可行体素集合定义为
\begin{equation}
\begin{aligned}
\mathbf{V}_{\mathrm{feas}}=& \{v_i\in\mathbf{V}\mid\ 
I_{\mathrm{occ}}(v_i)=0,\\
&\min_{b_j\in\mathbf{B}}d(\Omega(v_i),\Omega(b_j))\ge d_{\min}^{\mathrm{sfty}},h_{\min}\le z_i^v\le h_{\max}\}.
\end{aligned}
\end{equation}
其中，$d(\Omega(v_i),\Omega(b_j))$ 表示体素 $v_i$ 与建筑物 $b_j$ 占用区域之间的最小距离，$d_{\min}^{\mathrm{sfty}}$ 为安全距离阈值，$h_{\min}$ 和 $h_{\max}$ 分别为 UAV 的最小和最大飞行高度。

多无人机集合记为
$\mathbf{U}=\{u_1,u_2,\ldots,u_{|\mathbf{U}|}\}$。第 $n$ 架 UAV 的起始体素和目标体素分别为
$v_n^{\mathrm{start}}$ 和 $v_n^{\mathrm{end}}$。在离散时隙 $t\in\{0,1,\ldots,T\}$ 内，UAV 的位置由其所在体素表示为 $v_n(t)\in\mathbf{V}_{\mathrm{feas}}$，其轨迹表示为
\begin{equation}
\mathrm{Traj}_n=\{v_n(0),v_n(1),\ldots,v_n(T)\},
\end{equation}
其中 $v_n(0)=v_n^{\mathrm{start}}$，并期望在有限规划时域内到达 $v_n^{\mathrm{end}}$。

\subsection{Communication-Localization Fingerprint}

为统一表征城市环境中空间位置的通信和定位性能，本文提出通信-定位指纹（Communication-Localization Fingerprint, CLF）的概念。对于任意可行体素$v_i\in\mathbf{V}_{feas}$，其通信-定位指纹定义为：
\begin{equation}
F(v_i)=\{f_{comm}(v_i),f_{loc}(v_i)\},\quad v_i\in\mathbf{V}_{feas}.
\end{equation}
其中$F(\cdot)$表示体素到通信-定位特征空间的映射关系，$f_{comm}(v_i)$、$f_{loc}(v_i)$分别为通信指纹和定位指纹。

生成指纹信息前需要确定体素能够接入的地面基站集合。考虑建筑物遮挡对无线传播的影响，当体素中心与基站之间的连线不与任意建筑物占用区域相交，同时不超过基站的最大通信范围时，认为该基站与体素之间存在视距（Line-of-Sight, LoS）链路。

通信层面，本文采用体素位置处可获得的最大通信容量进行衡量。对于体素$v_i$与可见基站$m_k$，其信噪比可表示为$SNR_{i,k}=P_k\cdot G_{i,k}/\sigma_0^2$。其中，$P_k$为基站发射功率，$\sigma_0^2$为高斯白噪声功率，$G_{i,k}$为两者间的信道增益，具体表示为$G_{i,k}=\beta_0d(c_i^v,c_k^m)^{-\alpha}$，$\beta_0$为参考距离的信道增益，$\alpha$为路径损耗因子。根据香农公式，通信容量为$R_{i,k}=B\log(1+SNR_{i,k})$，其中$B$为系统带宽。进而得到体素$v_i$的通信性能为$C(v_i)=\max_{m_k\in\mathbf{M}_i}R_{i,k}$，$\mathbf{M}_i$为可行基站集合。

定位层面，采用TDOA（Time Difference of Arrival）的定位方法进行衡量。参考文献\cite{ref16}中给出的TDOA的定位误差下界公式进行计算得到体素$v_i$的定位性能为$L(v_i)=\sqrt{CRLB(v_i)}$

为了将通信性能与定位性能统一映射到相同量纲的特征空间，本文分别构建通信指纹与定位指纹。

对于通信指纹，若无可用基站，则将通信指纹设置为0，表示无法进行通信；反之，则将体素$v_i$的通信性能$C(v_i)$通过sigmoid函数映射为归一化的通信指纹，具体表示如下：
\begin{equation}
f_{comm}(v_i)=
\begin{cases}
\dfrac{
\begin{aligned}
&S\left(\alpha_C\left(C(v_i)-(C_{th}-\delta_C)\right)\right)\\[-0.2ex]
&-S\left(-\alpha_C(C_{th}-\delta_C)\right)
\end{aligned}}
{1-S\left(-\alpha_C(C_{th}-\delta_C)\right)}, & |\mathbf{M}_i|\ge 1;\\[3ex]
0, & \text{o.w.}
\end{cases}
\end{equation}
其中，$S(\cdot)$为sigmoid函数$S(x)=1/(1+e^{-x})$， $C_{th}$为通信性能门限，$\delta_C$、$\alpha_C$分别为通信侧门限偏移参数和函数斜率控制参数。

对于定位指纹，若其可用基站数量小于4时，认为该体素不满足三维TDOA定位所需的基站数量条件，将定位指纹设置为0；反之，则根据体素$v_i$的定位性能$L(v_i)$，通过sigmoid函数进行归一化，整合后的定位指纹定义如下：
\begin{equation}
f_{loc}(v_i)=
\begin{cases}
\dfrac{
\begin{aligned}
&S\left(\alpha_L\left(1/L(v_i)-1/(L_{th}+\delta_L)\right)\right)\\[-0.2ex]
&-S\left(-\alpha_L/(L_{th}+\delta_L)\right)
\end{aligned}}
{1-S\left(-\alpha_L/(L_{th}+\delta_L)\right)}, & |\mathbf{M}_i|\ge 4,\\[3ex]
0, & \text{o.w.}
\end{cases}
\end{equation}
其中， $L_{th}$为定位性能门限，$\delta_L$、$\alpha_L$分别为定位侧门限偏移参数和函数斜率控制参数。

\subsection{UAV Mobility Model}

在离散空间中，无人机在每个时隙只能从当前位置对应体素$v_n(t)$的邻接体素集合$\mathbf{V}_n^{adj}(t)$中选择下一体素，因此无人机轨迹是一串离散体素中心序列。为保证飞行速度约束，无人机仅能选取26邻域中的体素。对于体素边长为$\Delta r$的立方体，无人机单时隙位移距离取有限值：面邻接$\Delta r$、棱邻接$\sqrt{2}\Delta r$、角邻接$\sqrt{3}\Delta r$、悬停0。等效飞行速度$V_n(t)$被量化为有限集合$\{0,\Delta r/\Delta t,\sqrt{2}\Delta r/\Delta t,\sqrt{3}\Delta r/\Delta t\}$。

能耗方面，在时隙内将无人机近似为以恒定速度，沿相邻体素中心连线做匀速直线飞行，推进功率$P(V_n(t))$，采用现有旋翼无人机飞行功率模型 \cite{ref17}。则该时隙飞行能耗为$E_n^{fly}(t)=P(V_n(t))\cdot\Delta t$，无人机在规划时域内的总飞行能耗为$E_n^{fly}=\sum_{t=1}^{T}E_n^{fly}(t)$，$T$表示规划时域内的总时隙数。由于$V_n(t)$仅取有限离散值，每个动作对应的单步能耗可以离线预计算，用于后续奖励函数构建。

\subsection{Multi-UAV Planning Problem}

基于多无人机协同路径规划的最优化问题设计，主要是在满足约束条件的情况下，为每架无人机在每个时隙选择下一步可行体素，使多无人机系统在路径执行过程中获得较优的通信性能、定位性能和能耗表现。该问题表述如下：
\begin{equation}
\begin{aligned}
\textit{Problem:}\quad & \max_{\substack{\{v_n(t)\}\\ n=1,\cdots,|\mathbf{U}|\\ t=1,\cdots,T}}
\sum_{n=1}^{|\mathbf{U}|}\left\{\sum_{t=1}^{T}\left(w_Cf_{comm}(v_n(t))\right.\right.\\
&\left.\left.\quad +w_Lf_{loc}(v_n(t))\right)-w_EE_n^{fly}\right\}\\
\textit{s.t.}\quad & v_n(t)\in\mathbf{V}_{feas}\cap\mathbf{V}_n^{adj}(t-1),\\
&\quad \forall n=1,\cdots,|\mathbf{U}|,\ t=1,\cdots,T;\\
& v_n(t)\ne v_{n'}(t),\quad \forall n\ne n';\\
& v_n(0)=v_n^{start},\ v_n(T)=v_n^{end}.
\end{aligned}
\end{equation}
其中，$w_E$、$w_C$、$w_L$分别为能耗惩罚因子、通信需求因子和定位需求因子，其满足$w_C+w_L=1$；约束条件中，$v_n(t)\in\mathbf{V}_{feas}\cap\mathbf{V}_n^{adj}(t-1)$保证无人机始终在可行飞行区域内运动，同时保证当前时隙的无人机体素位置为上一时隙的邻域；$v_n(t)\ne v_{n'}(t)$避免多个无人机在同一时隙占据同一体素，从而保证飞行安全；$v_n(0)=v_n^{start},v_n(T)=v_n^{end}$保证无人机能够在有限的规划时域内顺利从起始点到达目标点。

虽然离线构建了通信—定位指纹，但执行阶段每架UAV只能观测到局部信息。因此该问题为部分可观测多智能体序贯决策问题。为此，下一节进一步将该问题建模为去中心化部分可观测马尔可夫决策过程（Dec-POMDP），并利用CTDE的多智能体强化学习框架进行求解。

\section{Proposed Algorithm}

在CTDE的框架下，使用MARL算法对最优化问题进行求解。本文采用多智能体近端策略优化（Multi-Agent Proximal Policy Optimization, MAPPO）作为Dec-POMDP的实现方法。定义如下：

\noindent\textbf{1. 状态空间建模}

在时隙$t$下，系统的全局状态由所有无人机的体素位置及其对应的通信-定位指纹构成，定义为：
\begin{equation}
s(t)=\{v_n(t),F(v_n(t))\}_{n=1,\cdots,|\mathbf{U}|}.
\end{equation}
在分布执行阶段，每架无人机的局部观测定义为：
\begin{equation}
o_n(t)=\{v_n(t),F(v_n(t)),\mathbf{V}_n^{adj}(t),F(\mathbf{V}_n^{adj}(t))\}.
\end{equation}
其中，$\mathbf{V}_n^{adj}(t)$、$F(\mathbf{V}_n^{adj}(t))$分别为体素$v_n(t)$的邻接体素集合及其指纹信息。

\noindent\textbf{2. 动作空间建模}

每个无人机在离散网格空间中移动，其动作定义为：
\begin{equation}
a_n(t)\in A_n=\{\mathbf{V}_n^{adj}(t)\}.
\end{equation}
动作本质为从当前体素的可行邻接体素集合中选择下一时隙的体素位置，即：$v_n(t)\xrightarrow{a_n(t)}v_n(t+1)$

\noindent\textbf{3. 奖励函数设计}

本文基于体素指纹信息构造团队平均的即时奖励函数。为保证各分量在相同量纲下可比，将通信、定位指纹直接采用其归一化取值，并对进展项与能耗项分别进行归一化处理。即时奖励函数定义为：
\begin{equation}
\begin{aligned}
r(t)=& \frac{1}{|\mathbf{U}|}\sum_{n}\Big\{X_n(t)\Big[
 w_C(f_{comm}(v_n(t))-1)\\
&+w_L(f_{loc}(v_n(t))-1)-w_E\widetilde{E}_n(t)-\delta_n^{coll}(t)\Big]\Big\}.
\end{aligned}
\end{equation}
其中，$X_n(t)$为到达状态指示变量，记为：
\begin{equation}
X_n(t)=
\begin{cases}
1, & v_n(t)\ne v_n^{end};\\
0, & \text{o.w.}
\end{cases}
\end{equation}
其表示无人机是否仍处于未完成任务状态。当$X_n(t)=0$时，认为该无人机已到达目标体素，其动作被锁定为悬停，不再累计通信指纹、定位指纹、进展奖励和飞行能耗。

$\widetilde{E}_n(t)$为归一化能耗惩罚，定义为：
\begin{equation}
\widetilde{E}_n(t)=\frac{E_n(t)}{E_{max}}.
\end{equation}
其中，$E_n(t)$为单步飞行能耗，$E_{max}$为最大单步能耗（对应角邻接移动）。

$\delta_n^{coll}(t)$为多无人机占据同一体素的惩罚项，定义为：
\begin{equation}
\delta_n^{coll}(t)=\lambda_{coll}I_{coll}(v_n(t)).
\end{equation}
其中，$\lambda_{coll}$为碰撞惩罚强度，$I_{coll}(v_n(t))$为碰撞指示函数，定义为：
\begin{equation}
I_{coll}(v_n(t))=
\begin{cases}
1, &\begin{aligned}[t]
&\exists n'\ne n,\ \text{s.t. }v_n(t)=v_{n'}(t)\ \text{or}\\
&v_n(t)=v_{n'}(t-1)\ \text{and}\ v_n(t-1)=v_{n'}(t);
\end{aligned}\\
0, & \text{o.w.}
\end{cases}
\end{equation}
碰撞指示函数刻画了同体素占据与相邻时隙换位穿越两类冲突。

综上，该奖励函数利用体素指纹的归一化特性实现通信与定位性能的统一度量，通过归一化能耗项与逐机碰撞项在统一量纲下约束飞行能耗与飞行安全。

基于上述状态空间、动作空间和奖励函数设计，MAPPO算法流程如下：

\begin{algorithm}[H]
\caption{CLF-Aided Multi-UAV Trajectory Planning via MAPPO}
\label{alg:clf_mappo}
\scriptsize
\begin{algorithmic}

\REQUIRE CLF map $F(v_i)$, feasible voxel set $\mathbf{V}_{\mathrm{feas}}$, 
start/goal voxels $v_n^{\mathrm{start}}$, $v_n^{\mathrm{end}}$, 
and MAPPO hyperparameters.

\ENSURE Shared policy network $\pi_{\theta}$.

\STATE Initialize shared actor $\theta$ and centralized critic $\phi$.

\REPEAT
    \STATE Reset UAV positions $v_n(0)\leftarrow v_n^{\mathrm{start}}$ and clear buffer $\mathcal{D}$.

    \FOR{$t=0,\ldots,T-1$}
        \STATE Each UAV constructs $o_n(t)$ from local CLF observations and selects a feasible action
        $a_n(t)\sim\pi_{\theta}(\cdot\mid o_n(t))$.
        \STATE Execute joint action, obtain team reward $r(t)$, observe $s(t+1)$, and store transition in $\mathcal{D}$.
        \IF{all UAVs reach their target voxels}
            \STATE \textbf{break}
        \ENDIF
    \ENDFOR

    \STATE Compute GAE advantages $\hat{A}(t)$ and returns $\hat{R}(t)$.
    \STATE Update actor $\theta$ using PPO clipped objective and update critic $\phi$ using value loss.

\UNTIL{episode $=Epi_{\max}$}

\end{algorithmic}
\end{algorithm}

\section{Experimental Evaluation}

本节通过仿真验证所提出的CLF辅助多无人机路径规划方法。仿真区域为 $600\mathrm{m}\times600\mathrm{m}\times150\mathrm{m}$ 的城市低空场景，其中包含15个建筑物和12个地面基站。三维空域通过体素离散为网格，UAV仅能在满足避障、高度（$h_{\min}=60\mathrm{m}$、$h_{\max}=120\mathrm{m}$）和安全距离（$d_{\min}^{\mathrm{safty}}=5\mathrm{m}$）约束的可行体素中移动。通信与定位门限 $C_{th}$、$L_{th}$ 不由人工设定，而是由可行体素上 $C(v_i)$ 与 $1/L(v_i)$ 的分布中位数自动标定，使通信指纹与定位指纹的过渡带各覆盖约 $30\%$--$70\%$ 的可行体素。主要MAPPO训练参数为：折扣因子 $\gamma=0.99$、GAE 参数 $\lambda=0.95$、裁剪系数 $\epsilon=0.2$、策略网络与评论家网络学习率分别为 $3\times10^{-3}$ 与 $5\times10^{-3}$、隐藏层宽度128、训练回合数2000。通信与定位权重保证$w_C+w_L=1$。

\begin{figure}[t]
\centering
\includegraphics[width=\columnwidth]{outputs/policy_comparison_3d.png}
\caption{三种规划器在同一城市场景中的三维路径：测地最短路参考策略（虚线）、无CLF奖励的MAPPO消融（点线）和所提CLF感知MAPPO（实线）。颜色区分无人机，方块与星号分别表示起点和终点。}
\label{fig:policy3d}
\end{figure}

\begin{figure}[t]
\centering
\includegraphics[width=0.9\columnwidth]{outputs/policy_comparison_bars.png}
\caption{三种规划器路径上的 CLF 质量（所有无人机均值）。三者均到达全部目标，因此 $f_{\mathrm{comm}}$、$f_{\mathrm{loc}}$ 及 CLF 奖励的差距完全来自是否优化 CLF，所提方法均显著占优。}
\label{fig:policy}
\end{figure}

图~\ref{fig:policy3d} 展示三种规划器的三维路径，图~\ref{fig:policy} 给出对应的 CLF 质量量化结果。对比的三种规划器为：测地最短路参考策略、不含CLF奖励的MAPPO消融，以及所提CLF感知MAPPO。三者均成功到达全部目标，从而排除了可达性差异的干扰。结果显示，所提方法在路径平均通信指纹与定位指纹上均明显高于两个基线，表明CLF奖励的引入使无人机主动绕开通信与定位性能退化区域，而非仅仅完成几何上的最短穿越。

\begin{figure}[t]
\centering
\includegraphics[width=\columnwidth]{outputs/algo_comparison.png}
\caption{三种代表性 on-policy AC 算法在 CLF 任务上的团队回报收敛曲线（阴影为种子间标准差）。MAPPO 收敛最快且回报最高；去除集中式评论家的 IPPO 明显最差，说明 CTDE 集中式评论家是关键；去除 PPO 裁剪的 MAA2C 居中，说明裁剪非核心因素。}
\label{fig:algo}
\end{figure}

图~\ref{fig:algo} 比较了三种代表性 on-policy AC 算法：MAPPO、MAA2C 与 IPPO 的团队回报收敛过程。结果表明：MAPPO 的团队回报收敛最快、最稳定，最终达到最高水平；去除集中式评论家的 IPPO 表现最差，而去除 PPO 裁剪的 MAA2C 与 MAPPO 接近。这说明所提框架的性能优势主要来自 CTDE 的集中式评论家，而非策略裁剪机制。在仅统计到达无人机的路径质量时，MAPPO 取得 $f_{\mathrm{comm}}=0.62$、$f_{\mathrm{loc}}=0.41$ 的均衡指纹质量。

\begin{figure}[t]
\centering
\includegraphics[width=\columnwidth]{outputs/resolution_training_return.png}
\caption{不同体素分辨率（$\Delta r=20/15/10\mathrm{m}$）下 MAPPO 的团队回报训练曲线（阴影为种子间标准差）。精细网格收敛较慢，但最终收敛至更高回报水平，说明更精细的体素网格能为规划器提供更高的 CLF 服务质量上界。}
\label{fig:resolution}
\end{figure}

图~\ref{fig:resolution} 给出体素分辨率对训练过程的影响。导航步长与奖励整形系数按 $1/\Delta r$ 缩放以保持物理尺度一致，三种分辨率使用相同的训练回合预算。结果表明，越精细的网格状态空间越大、收敛越慢，但在足够训练量下最终收敛至更高的团队回报，说明更精细的体素网格能为规划器提供更高的 CLF 可达服务质量上界。这揭示了环境表达精度与训练复杂度之间的折中：精细网格提升了理论可达的服务质量，但需要更多样本才能收敛。

\begin{figure}[t]
\centering
\includegraphics[width=0.9\columnwidth]{outputs/pareto_clf.png}
\caption{通信与定位权重 $w_L$ 对路径平均指纹的影响。刻画城市遮挡环境下通信与定位权衡。}
\label{fig:pareto}
\end{figure}

图~\ref{fig:pareto} 展示了通信与定位权重对路径规划结果的影响。随着 $w_L$ 增大，路径上的平均定位指纹 $f_{\mathrm{loc}}(v_i)$ 由0.37单调提升至0.64，而平均通信指纹 $f_{\mathrm{comm}}(v_i)$ 相应由0.69下降至0.54，且各权重下三架无人机均成功到达目标。该结果说明城市低空遮挡环境中通信质量和定位可靠性存在明显空间权衡，所提出CLF框架能够通过权重调节适应不同通定服务需求。

\section{Conclusion}

本文研究了城市低空环境下面向通信与定位双需求的多无人机协同路径规划问题，提出了一种基于通信-定位指纹（CLF）的协同规划框架。通过结合通信-定位指纹地图与多智能体强化学习方法，实现了复杂多约束条件下的协同路径规划。本文首次将通信与定位性能统一映射为空间指纹并引入多无人机路径规划过程，从而避免了在线重复评估通信与定位性能。仿真结果表明，所提出的方法能够有效提高多无人机任务执行过程中的服务连续性和路径规划可靠性。

\begin{thebibliography}{17}
\bibliographystyle{IEEEtran}

\bibitem{ref1} 
Integrated sensing and communications: Toward dualfunctional wireless networks for 6G and beyond
\bibitem{ref2} 
UAV-assisted data collection for Internet of Things: A survey
\bibitem{ref3} 
Survey of important issues in UAV communication networks
\bibitem{ref4} 
Survey on unmanned aerial vehicle networks for civil applications: A communications viewpoint
\bibitem{ref5} 
A survey on millimeter-wave beamforming enabled UAV communications and networking
\bibitem{ref6} 
Wireless communications with unmanned aerial vehicles: Opportunities and challenges
\bibitem{ref7} 
3D Urban UAV Relay Placement: Linear Complexity Algorithm and Analysis
\bibitem{ref8} 
Optimizing UAV Deployment and Access Control for Integrated Localization and Communication
\bibitem{ref9} 
A Hierarchical MADRL Framework for Joint Resource Allocation and Trajectory Design in UAV Enabled Multi-Static ISAC Networks
\bibitem{ref10} 
Transformer based collaborative reinforcement learning for fluid antenna system (FAS)-enabled 3D UAV positioning
\bibitem{ref11} 
GeoAgg-HSAC: an RL-based Framework for Trajectory and Resource Optimization in Mountainous UAV Integrated Localization and Communication Networks
\bibitem{ref12} 
MARL-Based UAV Trajectory and Beamforming Optimization for ISAC System
\bibitem{ref13} 
UAV-aided RF mapping for sensing and connectivity in wireless networks
\bibitem{ref14} 
UAV-aided wireless node localization using hybrid radio channel models
\bibitem{ref15} 
Joint 3D positioning and power allocation for UAV relay aided by geographic information.
\bibitem{ref16} 
Achievable Localization Accuracy of the Positioning Reference Signal of 3GPP LTE.
\bibitem{ref17} 
Y. Zeng, J. Xu, and R. Zhang, ``Energy minimization for wireless communication with rotary-wing UAV,'' IEEE Trans. Wireless Commun., vol. 18, no. 4, pp. 2329--2345, Apr. 2019.

\end{thebibliography}

\end{document}
